\documentclass[12pt]{article}
\usepackage[utf8]{inputenc}
\usepackage[T1]{fontenc}
\usepackage{hyperref}
\title{Xiri Language Resources for K-12 Curriculum Development}
\author{Compiled by Claude Sonnet 5.5 (claude-sonnet-5-5) for Andr\'as Kornai}
\date{}
\begin{document}
\maketitle

\section{Introduction}
Xiri (ISO 639-3 xii; Glottocode xiri1242) is an endangered Khoe-Kwadi language of South Africa. The reference description says it is thought to be used as a first language only by the elderly, with about 9 speakers, and that it is not known to be taught in schools. My searches found very little material on Xiri itself. Sources I saw treat Xiri as one of several names (with Korana and Griqua) for the Khoemana or South Khoekhoe language, and one linguistic paper questions the spelling ``Xiri''. Speaker figures differ between sources (ASJP lists about 187 speakers, dated 2008), so none of them should be treated as settled.

\section{Primary Educational Resources}
I found no primers, textbooks, or school programmes for Xiri. Any curriculum would have to be built from scratch, with community involvement.

\section{Linguistic and Reference Materials}
\begin{itemize}
\item Khoemana (Wikipedia): overview of the closely related Khoemana language, including its sound system compared with Khoekhoe. \url{https://en.wikipedia.org/wiki/Khoemana}
\item ASJP database entry for Xiri: a word list and basic data. \url{https://asjp.clld.org/languages/XIRI}
\item Channumerals entry, ``Xiri, Western Cape, South Africa'': numeral data. \url{https://lingweb.eva.mpg.de/channumerals/Xiri.htm}
\item G\"uldemann, conference paper on Western Ui (Khoisan) languages. \url{https://www.eva.mpg.de/fileadmin/content_files/linguistics/conferences/2015-speaking-of-Khoisan/P5_Linguistics_Gueldemann_Western__Ui.pdf}
\item Article in \emph{Studies in African Linguistics} that discusses the spelling of the name Xiri. \url{https://journals.flvc.org/sal/article/download/118040/116005/172172}
\end{itemize}

\section{Religious and Community Resources}
I found no scripture translations or community organisations specific to Xiri.

\section{Digital and Audio Resources}
I found no audio, video, or corpus resources specific to Xiri.

\section{Curriculum Development Context}
There is almost no documented material for Xiri. Closely related Khoekhoe-group materials might serve as a starting point, but I did not verify any specific ones. A realistic first step is documentation work with the remaining elderly speakers, together with the sources above.
\end{document}